\section{Methods for modelling metaporous surfaces and interfaces}
\subsection{A word on periodic structures}
A large literature that shows the enhancement of absorption properties by periodicizing porous layers with scatterers was discussed in the introduction. In this section, a presentation of the fundamental notions surrounding periodic structures is proposed. In wave physics, periodic structures have been shown to introduce interesting effects towards applications for wave control and constitute the cornerstone of acoustic metamaterials. The periodicity can be imposed in acoustic systems using the so-called Bloch-Floquet theorem. This formalism has been pioneered by Felix Bloch and Léon Brillouin among others \cite{bloch1929,brillouin1946}. Throughout this manuscript, we consider a periodic cell along only the longitudinal direction, that 1D periodicity which we denote $d$. Taking an infinitely repeating lattice or a crystal, the description of the wave propagation in only one of the so-called unit cell, one spatial period of the total structure, is enough to describe the whole structure. The layer of infinite extent is thus limited to a single unit-cell. A wavefunction $\psi(\mbf x)$ describing the propagation should thus depend on a field with the same spatial periodicity as the lattice. To do so, the wavefunction is projected onto an infinite modal basis, that is Bloch modes,
\begin{equation}
​	\psi_q (\mbf x) = e^{\im k_q x_1}.
\end{equation}
The wavefunction at the left boundary is thus related to that at the right boundary modulated by the spatial periodicity and some translation vector, that is $\psi(x + q d) = e^{\im \mbf k} \psi(x)$. 

The possibility of reduction of the spatial considerations of the problem also leads to a reduction of the reciprocal wavenumber space. In periodic structure, we often restrict this reciprocal space to the first Brillouin zone, that is to wavenumbers restricted to $k_{1q} \in [0, \pi/d]$ when looking at the band structure. The consequences of periodic media are directly observable in the band structure: bands are usually folded when the wavenumber reaches this limit of the reciprocal space. This folding occurs at the Bragg frequency corresponds to a relation between the wavelength $\lambda$ of the excitation signal and the periodicity of the crystal lattices, $d \sin(\theta) = q \lambda$, with $\theta$ the angle of incidence. This condition corresponds to destructive interferences: this phenomenon is a cause of the appearance of the famous Bragg bandgaps; in these bands of frequency, waves are unable to propagate. Variations of these phenomenon have been investigated in acoustic metamaterials: e.g. the addition of periodic resonators leads to hybridization stopbands. \comment{ref these}

\subsection{Methods to compute the cattering of waves by periodic structures} \comment{manque des refs partout }
Analytical methods for the evaluation of the wave propagation in periodic structures often require tedious derivations, as the frequency regimes of interest is usually high enough to require the description of the multiple scattering happening in the structure. The scatterer of the unit cell onto which the incident wave impinges acts as a secondary source, which scattered field is successively scattered by each neighboring obstacles and so on. \acrlong{mst} is a theoretical framework that has been developed to account for these phenomena and was used to compute the acoustic response of metaporous structures with rigid inclusions \cite{groby2011} and metaporoelastic surfaces with solid and shell inclusions \cite{weisser2016}. Both configurations have been tuned to show a perfect absorption peak in the frequency spectrum. 

The remaining papers discussing how incident pressure fields are affected and scattered by metaporous structures use \acrshort{fem}, for the most part through the interface of Comsol Multiphysics. This very versatile and powerful method has been used for decades in all fields of physics and thus won't be detailed in this thesis. It is recommended to the reader as an introduction on the basis of \acrshort{fem}, the textbooks from {Atalla} and {Sgard}\cite{atalla2015} related to acoustics and vibration problems and that of Dhatt \etal \cite{dhatt2005} and Solin \etal \cite{solinHigherorderFiniteElement2004} dealing specifically with higher-order finite elements problems. This method allows for the discretization of the geometry via a mesh in order to solve the appropriate motion equations. This has allowed to use rather complex geometries, e.g spirals \cite{zhong2019} and labyrinthine paths \cite{liuAcousticLabyrinthinePorous2021} as scatterers embedded in porous layers as discussed in the state of the art. Additionally, Gaborit \etal \cite{gaborit2018} have developed a method in order to couple \acrshort{fem} with Bloch modes expansion of analytical fields described by \acrshort{tmm}. They have used it to compute the scattering coefficients of structure that include a periodic cell sandwiched by homogeneous layers with an embedded inclusion.

\subsection{Finite elements}
\comment{y a du taf ici :/ manque en gros 2 pages environ}
\begin{itemize}
    \item reprendre sur comment arriver à NME eq.22: formulation eq. faible + maillage + interpolation polynomiale (higher order) aux nodes du mesh
    \item sert à etoffer la partie du chapitre 4: écrire la recette permettant d'y arriver.
    \item cas du fluide et montrer les différentes variational formulations
\end{itemize}

\section{Dispersion relations and band structure in periodic lattices}
\label{sec:disp_perioic_chap1}
\subsection{Methods for non-radiative systems}
One of the most popular methods for the computation of band structures in periodic systems made of classical lossless materials are the \acrfull{pwe} which is semi-analytical. It relies on a Fourier transform of the motion equations, which allows an expression of the fields expanded onto Bloch modes, truncated for numerical applications. This method, while easy to implement presents important limitations due to the convergence of Fourier series and the matching of constituent materials for the lattice: materials should preferably of the same nature for the calculation to go well\cite{jimenez2021}. Generalized as the \acrfull{epwe}, the method also allows to obtain the evanescent Bloch modes in Ref.~\cite{schalcher2023}, allowing for complex dispersion diagram investigation and to account for evanescent modes in finite 2D periodic lattices. 

The complexity of the system prohibits the use of such semi-analytical methods. For more complex geometries, we must turn to alternative methods. As the 1D discretization used by the \acrshort{scm} and \acrshort{safe} and the plane wave formalism used in \acrshort{tmm} and the \acrshort{gm} is not capable of providing a description of the geometry of inhomogeneous layers, in our case, the effects of the scatterer in the unit cell on wave propagation. The thesis of Magliacano\cite{magliacano2020a} was focused in the modeling of a closed system, of an equivalent fluid unit cell with an embedded rigid inclusion. The finite elements formalism was entitled with the so-called shift cell technique, that implements the Bloch-Floquet phase shift at the boundaries inside of the constitutive variational weak formulation. 

% The \acrshort{wfe} as well as classical \acrshort{fem} have been applied to analogous systems. Salam Claro has developed a \acrshort{wfe} method in his thesis \cite{salam_claro_phd} to model elasto-acoustic waveguides in 3D geometries. Their modularity and versatility lead them to be good candidates for our configurations of interest. However, these methods have lacked developments to model open domains allowing the waves propagating in the structure to radiate out of the structure.  

\subsection{Unpromising adaptation of existing methods}
The complexity of the system of interest in this thesis as presented in Section\ref{jsp} of the Introduction prohibits the use of such semi-analytical methods. For more complex geometries, we must turn to alternative methods: as the 1D discretization used by the \acrshort{scm} and \acrshort{safe} and the plane wave formalism used in \acrshort{tmm} and the \acrshort{gm} is not capable of providing a description of the geometry of inhomogeneous layers, in our case, the effects of the scatterer in the unit cell on wave propagation. 

Multiple scattering could also be considered. However, it has been successfully used to describe the acoustic scattering of periodically arranged rigid inclusions in equivalent fluid unit cells \cite{groby2011} as well as solid inclusions in a poroelastic unit cell\cite{weisser2016}. However, the so-called Schlömilch series that appear in the calculation, from the effect of the scattering of the neighboring cylindrical inclusions appear to be not so manipulable when considering complex longitudinal wavenumbers. Its expression is  
\begin{equation}
    \sigma_q^X=\sum_{l=1}^\infty H_q^{(1)}(k^Xld)\left((-1)^qe^{ik_1^Xld}+e^{-ik_1^Xld}\right),
\end{equation}
and indeed, separating the $k_1^X$ real and imaginary inevitably leads to at least one growing exponential term in the series, leading to it being unable to converge, up to our study. 

We face challenges when trying to model this structure using the tools classically used in the literature, that is COMSOL Multiphysics. Although there are built-in tools in the software for "\emph{Eigenfrequency}" studies, the library aims at solving complex frequencies of the system, given a set of parameters for the wavenumber. As it has been discussed earlier in the chapter, the goal of this work is to solve for complex wavenumbers and real frequencies. Even trying a manual implementation of the model in the "\emph{Weak Form PDE}" interface, the tendency of this software to function as a black box introduces a new degree of complexity to this strategy. While \acrshort{fem} remains the soundest strategy for the modeling of the complex inhomogeneous geometry of the unit cell of metaporous surface, we propose to adapt in-house existing frameworks in order to develop numerical tools in this thesis.  

\section{Conclusion and prospects of methods for the thesis}

In the case of homogeneous multilayer systems with surrounding fluid loading, we have seen that semi-analytical have long been established in the literature in order to compute their scattering properties. For dispersion problems however, these configurations usually require some sort of numerical discretization to treat the geometry; a wide range of methods allowing for the direct formulation of an eigenvalue problem and a computation of modes was shown in Section \ref{sec:chap1_guidedwaves} with first semi-analytical tools applicable to simple cases (e.g. an equivalent fluid layer with surrounding air), and then various numerical methods like SCM or SAFE as the two main protagonist.

When going to periodic systems, it appeared that the way to tackle the problem is quite different. Indeed, the resulting complex geometry (due to the obstacles) leads to much different tools. In the case of scattering properties computation, semi-analytical tools are still available. We have seen that \acrshort{fem} seems to provide promising possibilities for the modelling of complex periodic cores. However, the considerations of open systems leads to question regarding the modeling of the half-spaces surrounding the structure. Two strategies compete in the literature: either a fluid cavity is added to the geometry and terminated by some sort of absorbing layer at the end; or the harmonic fluid radiation is locally accounted for as an analytical interface condition of the structure. \comment{parler de couplage avec plaque et sondwich}